\chapter{Software e Firmware}
\label{chapter:software}

\begin{introduction}
Este capítulo descreve o software e o firmware desenvolvidos para a mesa sísmica. A primeira secção apresenta o modo de operação principal - movimento sinusoidal bi-axial com parâmetros configuráveis via interface web embarcada. A segunda secção descreve o pipeline Python de conversão de dados sísmicos, incluindo as estratégias de adaptação do perfil e as razões pelas quais nenhuma delas resolve as limitações mecânicas fundamentais.
\end{introduction}



\section{Modo de Operação Principal: Frequência Variável com Web UI}
\label{sec:modo-principal}

O modo de operação principal da mesa é o movimento sinusoidal bi-axial com parâmetros configuráveis em tempo real via interface web embarcada no ESP32. Este é o único modo de operação validado experimentalmente e que cumpre os requisitos funcionais definidos.


\subsection{Arquitetura do Firmware Web}
\label{subsec:firmware-web}

O firmware implementado no ESP32 usa exclusivamente bibliotecas do núcleo Arduino para ESP32: \texttt{WiFi.h} para a configuração da rede e \texttt{WebServer.h} para o servidor HTTP - sem dependências externas. O ESP32 arranca em modo Access Point com as seguintes configurações:

\begin{itemize}
    \item \textbf{SSID:} \texttt{MesaSismica}
    \item \textbf{Password:} \texttt{12345678}
    \item \textbf{IP:} \texttt{192.168.4.1}
\end{itemize}

Qualquer dispositivo com browser (computador, tablet ou smartphone) pode ligar-se à rede \texttt{MesaSismica} e aceder à interface de controlo em \texttt{http://192.168.4.1}, sem necessidade de ligação USB ou computador externo durante o ensaio.


\subsection{Interface Web - Parâmetros Configuráveis}
\label{subsec:interface-web}

A interface web serve uma página HTML com os seguintes controlos:

\begin{itemize}
    \item \textbf{Frequência (Hz):} slider para definir a frequência da oscilação sinusoidal, entre 0,5\,Hz e 10\,Hz.
    \item \textbf{Amplitude (steps):} slider para definir o número máximo de passos que a mesa se pode deslocar.
    \item \textbf{Intervalo (ms):} slider para definir a janela de atualização da velocidade.
    \item \textbf{Botões Start / Stop:} iniciam e interrompem o movimento sem necessidade de reconectar ou reprogramar o ESP32.
\end{itemize}

Quando o utilizador submete novos parâmetros, o servidor HTTP recebe os valores via pedido HTTP GET, atualiza as variáveis globais e aplica-as imediatamente ao ciclo de controlo, sem interromper o movimento em curso.


\subsection{Geração do Movimento Sinusoidal}
\label{subsec:movimento-sinusoidal}

O movimento da plataforma é gerado por duas funções sinusoidais desfasadas de 90°, uma para cada eixo:

\begin{equation}
v_{\text{NS}}(t) = A \cdot \sin(2\pi f t)
\label{eq:vel-ns}
\end{equation}

\begin{equation}
v_{\text{EW}}(t) = A \cdot \sin\!\left(2\pi f t + \frac{\pi}{2}\right) = A \cdot \cos(2\pi f t)
\label{eq:vel-ew}
\end{equation}

onde $A$ é a amplitude configurada (steps/s), $f$ é a frequência configurada (Hz) e $t$ é o tempo decorrido desde o início do movimento (s). A desfasagem de 90° entre NS e EW produz um percurso elíptico (ou circular, quando a amplitude é igual nos dois eixos) da plataforma. A velocidade instantânea de cada motor é calculada a cada janela temporal e passada a \texttt{setSpeed()}, com \texttt{runSpeed()}~\cite{AccelStepper} chamado continuamente no \texttt{loop()} para distribuir os passos ao longo da janela.

O Código~\ref{lbl:firmware-web} ilustra a estrutura do ciclo de controlo do firmware Web:

\begin{listing}[h]
\begin{minted}{cpp}
// Variáveis globais atualizadas pelo servidor HTTP
float freq_hz    = 1.0;   // Hz
float amplitude  = 200.0; // steps/s de pico
float janela_ms  = 20.0;  // ms
bool  em_marcha  = false;

unsigned long ultimoUpdate = 0;
float t_decorrido = 0.0;  // segundos desde Start

void loop() {
    server.handleClient();  // Processa pedidos HTTP

    if (em_marcha && millis() - ultimoUpdate >= (unsigned long)janela_ms) {
        ultimoUpdate = millis();
        t_decorrido += janela_ms / 1000.0;

        float velNS = amplitude * sin(2.0 * PI * freq_hz * t_decorrido);
        float velEW = amplitude * cos(2.0 * PI * freq_hz * t_decorrido);

        motorN.setSpeed( velNS);
        motorS.setSpeed(-velNS);
        motorE.setSpeed( velEW);
        motorW.setSpeed(-velEW);
    }

    motorN.runSpeed();
    motorS.runSpeed();
    motorE.runSpeed();
    motorW.runSpeed();
}
\end{minted}
\caption{Ciclo de controlo do firmware Web: movimento sinusoidal bi-axial com parâmetros configuráveis.}
\label{lbl:firmware-web}
\end{listing}


\subsection{Validação Experimental}
\label{subsec:validacao-web}

O modo de frequência variável foi validado experimentalmente: a plataforma executa movimento bi-axial simultâneo nos dois eixos, com percurso elíptico observável a olho nu. Os parâmetros foram alterados em tempo real via browser durante o ensaio, sem qualquer reconexão USB. Os resultados completos são apresentados no Capítulo~\ref{chapter:validation}.



\section{Pipeline de Dados Sísmicos: Análise e Validação das Limitações}
\label{sec:pipeline-dados}

O pipeline de dados sísmicos é o elo entre os dados históricos do registo de El Centro e o movimento físico da plataforma. O processo tem três etapas: obtenção e preparação dos dados de deslocamento em formato CSV, conversão para sequências de passos de motor por um script Python e geração dos arrays C++ a embutir no firmware. Este pipeline corre em modo \textit{offline} e foi desenvolvido para validação do pipeline de conversão e análise das limitações mecânicas do sistema - não constitui o modo de operação da mesa.


\subsection{Fonte de Dados - Registo de El Centro}
\label{subsec:fonte-dados}

O registo sísmico de El Centro de 1940, formalmente designado por \textit{Imperial Valley Earthquake} de 18 de maio de 1940~\cite{ElCentro1940}, é um dos registos acelerométricos mais utilizados e documentados na literatura de engenharia sísmica. Com uma magnitude de aproximadamente 6,9 na escala de Richter, o sismo afetou a região do Imperial Valley, na Califórnia (EUA). O registo obtido na estação de El Centro (Imperial Valley Irrigation District) tornou-se referência na literatura de engenharia sísmica experimental, por ter sido um dos primeiros registos de forte movimento do solo com qualidade instrumental adequada e por abranger uma gama de frequências relevante para estruturas de edifícios correntes.

O registo é composto por duas componentes horizontais ortogonais: Norte-Sul (NS) e Este-Oeste (EW), cada uma disponível em séries temporais de aceleração, velocidade e deslocamento do solo, com um intervalo de amostragem de 20\,ms. Para o presente projeto, foram utilizados os dados de deslocamento (em metros), por serem a grandeza diretamente relevante para o controlo de posição da plataforma. O ficheiro NS apresenta 2687 amostras e o ficheiro EW apresenta 2673, o que corresponde a aproximadamente 53 segundos de registo sísmico.

Os dados sísmicos utilizados foram fornecidos pelo orientador do projeto, Professor Vítor Silva, sob a forma de um ficheiro Microsoft Excel (\texttt{.xlsx}), denominado \texttt{elcentro\_eq\_disp.xlsx}~\cite{PEERDatabase}. O ficheiro continha duas folhas de cálculo, uma para cada componente horizontal do registo: Norte-Sul (NS) e Este-Oeste (EW). Cada folha apresentava três colunas: tempo (em segundos), aceleração (em m/s\textsuperscript{2}) e deslocamento (em metros), com um intervalo de amostragem de 20\,ms. Os dados de deslocamento foram utilizados diretamente pelo script Python, sem necessidade de integração numérica adicional, uma vez que o ficheiro já continha os valores de deslocamento processados.

\subsection{Conversão de Deslocamento para Passos de Motor (Python)}
\label{subsec:conversao-python}

A conversão dos dados de deslocamento em passos de motor baseia-se no princípio de que cada janela temporal de 20\,ms corresponde a um incremento de deslocamento linear da plataforma, que deve ser reproduzido pelo motor num número inteiro de passos. O incremento de deslocamento entre duas amostras consecutivas $i-1$ e $i$ é dado por:

\begin{equation}
\Delta d_i = d_i - d_{i-1} \quad [\text{m}]
\label{eq:delta-deslocamento}
\end{equation}

Este incremento é então convertido para o número de passos de motor correspondente, a partir do \textit{lead} do fuso, do ângulo de passo do motor e do modo de \textit{microstepping} configurado:

\begin{equation}
\text{steps}_i = \left\lfloor f_{\text{ajuste}} \cdot \frac{\Delta d_i \times 1000}{\text{lead}} \cdot \frac{360^\circ}{\theta_{\text{passo}}} \cdot M_{\text{micro}} \right\rfloor
\label{eq:steps}
\end{equation}

onde $\Delta d_i \times 1000$ converte metros em milímetros, $\text{lead} = 8\,\text{mm/rev}$, $\theta_{\text{passo}} = 1{,}8^\circ$ (resultando em 200 passos/rev), $M_{\text{micro}}$ é o fator de \textit{microstepping} e $f_{\text{ajuste}}$ é um fator de escala que permite ajustar a amplitude do movimento à gama de deslocamentos fisicamente praticável pela plataforma.

O Código~\ref{lbl:pipeline-python} apresenta a função Python responsável por este cálculo, aplicada em simultâneo às componentes NS e EW do registo sísmico.

\begin{listing}[h]
\begin{minted}{python}
def steps_calc(data1, data2, max_steps_per_window=50, max_position=2500,
               microstepping=2, lead=8, step_angle=1.8, subsample=2):
    # Numero de passos por revolucao em full-step (200 para NEMA17)
    steps_per_revolution = 360 / step_angle

    steps_N, steps_E, steps_S, steps_W = [], [], [], []
    pos_NS = 0   # Posicao acumulada eixo NS (para controlo de limites)
    pos_EW = 0   # Posicao acumulada eixo EW (para controlo de limites)

    # Subsampling: percorre apenas 1 em cada 'subsample' amostras
    # Equivale a reproduzir o sismo em camara lenta (fator subsample)
    indices_ns = range(1, len(data1.displacement), subsample)
    indices_ew = range(1, len(data2.displacement), subsample)

    prev_ns = data1.displacement.iloc[0]
    for i in indices_ns:
        curr = data1.displacement.iloc[i]
        # Incremento de deslocamento entre amostras, em milimetros
        delta_mm = (curr - prev_ns) * 1000
        prev_ns = curr

        # Conversao de deslocamento para passos de motor
        s = round((delta_mm / lead) * steps_per_revolution * microstepping)

        # Clamping: limita a velocidade maxima por janela
        s = max(-max_steps_per_window, min(max_steps_per_window, s))

        # Controlo de posicao: impede que a plataforma ultrapasse os limites fisicos
        nova_pos = pos_NS + s
        if nova_pos > max_position:
            s = max_position - pos_NS
        elif nova_pos < -max_position:
            s = -max_position - pos_NS

        pos_NS += s
        steps_N.append(s)
        steps_S.append(-s)  # Motor S e oposto ao motor N (contrafase)

    prev_ew = data2.displacement.iloc[0]

    # O ciclo para a componente EW e identico, com data2 e pos_EW

    # Trunca os arrays ao menor comprimento (NS e EW podem ter tamanhos diferentes)
    min_len = min(len(steps_N), len(steps_E))
    steps_N = steps_N[:min_len]
    steps_S = steps_S[:min_len]
    steps_E = steps_E[:min_len]
    steps_W = steps_W[:min_len]

    # Informacao de debug: intervalo real do firmware apos subsampling
    print(f"// totalPontos = {min_len}")
    print(f"// Pos NS final: {pos_NS} | Pos EW final: {pos_EW}")
    print(f"// Intervalo firmware: {20 * subsample}ms")

    return steps_N, steps_E, steps_S, steps_W
\end{minted}
\caption{Função Python de conversão de deslocamento sísmico em passos de motor.}
\label{lbl:pipeline-python}
\end{listing}

\subsection{Parâmetros - Lead do Fuso, Microstepping, Steps por Revolução}
\label{subsec:parametros-conversao}

A Tabela~\ref{tab:parametros-conversao} resume os parâmetros fixos e configuráveis do pipeline de conversão, com os respetivos valores e significado físico.

\begin{table}[h]
\centering
\caption{Parâmetros do pipeline de conversão de deslocamento sísmico em passos de motor.}
\label{tab:parametros-conversao}
\begin{tabular}{lll}
\hline
\textbf{Parâmetro} & \textbf{Valor} & \textbf{Descrição} \\
\hline
Lead do fuso & 8\,mm/rev & Deslocamento linear por revolução completa \\
Ângulo de passo (NEMA17) & 1{,}8\textdegree{} & Passo nominal do motor em \textit{full-step} \\
Passos por revolução (\textit{full-step}) & 200 & $= 360^\circ / 1{,}8^\circ$ \\
Microstepping configurado & $1/2$ & Jumper M0 na CNC Shield \\
Resolução linear por passo & $0{,}02$\,mm & $= 8\,\text{mm} / (200 \times 2)$ \\
Fator de ajuste ($f_{\text{ajuste}}$) & variável & Escala de amplitude; $f < 1$ atenua o movimento \\
Clamping (\texttt{max\_steps\_per\_window}) & configurável & Limita passos máximos por janela \\
Posição máxima (\texttt{max\_position}) & configurável & Limite físico acumulado da plataforma (passos) \\
Subsampling & configurável & Executa 1 em cada $N$ amostras \\
Intervalo temporal base & 20\,ms & Duração de cada janela de dados originais \\
Frequência de atualização base & 50\,Hz & $= 1 / 0{,}02\,\text{s}$ \\
Amostras NS (El Centro) & 2687 & Duração $\approx$ 53,7\,s \\
Amostras EW (El Centro) & 2673 & Duração $\approx$ 53,5\,s \\
\hline
\end{tabular}
\end{table}


\subsection{Estratégias de Adaptação do Perfil Sísmico}
\label{subsec:estrategias-adaptacao}

Para tentar viabilizar a execução do perfil sísmico sem perda de passo, foram analisadas três estratégias de adaptação do sinal. Nenhuma resolve o problema de fundo - apenas muda a forma como o perfil é distorcido.

\textbf{Fator de ajuste ($f_{\text{ajuste}} < 1$):} variável de escala multiplicativa aplicada a todos os $\Delta\text{steps}$ (Equação~\ref{eq:steps}). Com $f_{\text{ajuste}} < 1$, a amplitude global do movimento é atenuada, mas a relação temporal entre amostras mantém-se, preservando o perfil dinâmico relativo do sismo. A desvantagem é que, com um $f_{\text{ajuste}}$ suficientemente baixo para evitar perda de passo, o deslocamento real da plataforma fica a uma velocidade estremamente baixa - imperceptível e sem significado para ensaio de modelos estruturais.

\textbf{Clamping (\texttt{max\_steps\_per\_window}):} limitação do valor absoluto de steps por janela a um máximo predefinido. Esta abordagem preserva os picos de baixa amplitude mas trunca os picos mais intensos, distorcendo precisamente o conteúdo em frequência e amplitude que é mais relevante para a resposta estrutural.

\textbf{Subsampling:} redução da frequência de atualização, executando apenas 1 em cada $N$ amostras consecutivas. Equivale a reproduzir o sismo em câmara lenta: com $N = 4$, o perfil é executado a um quarto da velocidade real. Os motores conseguem acompanhar - mas um sismo reproduzido 4$\times$ mais lento não é um sismo: as frequências de excitação baixam para fora da gama de ressonância típica de estruturas, e o ensaio perde validade sísmica.

Apesar destas estratégias, a análise experimental revelou que o sistema, nas condições atuais, não consegue reproduzir o registo sísmico de El Centro à escala temporal real. As razões mecânicas são analisadas em detalhe no Capítulo~\ref{chapter:validation}.
